% ---------------------------------------------------------------------------
% Appendix B: Detailed methods (empirical).
% Math lives in the theory appendix; this part gives data, settings and
% numerics.  Plan and sources: notes/appendix_outline_plan.md.
% ---------------------------------------------------------------------------

\section{Detailed methods}
\label{sec:app-methods}

% ---------------------------------------------------------------------------
\subsection{Natural-image data and covariance}
\label{sec:app-methods-images}
% Used by: Fig.~\ref{fig:phenomena} (bottom), Fig.~\ref{fig:pixel},
% Fig.~\ref{fig:feature}, App.~\ref{sec:app-ext-pixel}, \ref{sec:app-ext-gauge}.
% Content: van Hateren database; 100x100 patches; log(1+L)/log(65536),
% 8-bit quantization, rescale to [0,1]; 2,000-patch training pool (one per
% image) and 20,000-patch population pool (ten per image, disjoint source
% images); centered population covariance Sigma and its eigendecomposition;
% disk teacher (radius 0.3, S = 1736.34).  Currently stated inside
% App.~\ref{sec:app-vanhateren-landscape-methods}; move the shared part here
% and keep only the landscape-specific part there.
% Optional small panel: spectrum s_k and teacher power b_k^2
% (tables/vanhateren_spectral_teacher_profile.csv).
\todo{Write: images, preprocessing, covariance, disk teacher.}

% ---------------------------------------------------------------------------
\subsection{Toy teacher--student example}
\label{sec:app-methods-toy}
% Used by: Fig.~\ref{fig:phenomena} (bottom row).
% SOURCE: Closed-loop-visual-insilico/scripts/accentuation_theory/
%   exp2_accentuation.py (commit 09202ed; first run 2026-03-31); outputs in
%   $STORE_DIR/DL_Projects/AdvExampleLinearRegr/exp2/ (exp2_weights,
%   exp2_natural_scatter, exp2_accentuated_images_start*_highcontr,
%   exp2_cross_eval).  Companion: exp1_adversarial_weight_construction.py,
%   exp1_supp_eigenmode_gallery.py (eigenmode gallery; possible B.1 panel).
%   Precursor notebook: notebooks/20250528_toy_model_adv_generation.ipynb
%   (OLS/Ridge/Lasso fits of the disk teacher, mostly FFHQ).
% Settings (verified: target values in Fig. 1B match the script output):
%   van Hateren .iml (raw, 1024x1536), 500 images chosen with seed 42,
%   12,000 random 100x100 patches, per-image min-max normalization, offset
%   -0.15 (NB: differs from the log-luminance preprocessing of B.1/B.4);
%   Sigma from SVD of the centered patch matrix.
%   w* = disk mask, radius 0.3 on [-1,1]^2.  Students are CONSTRUCTED, not
%   fitted: w_1 = w* + 1 x unit direction in eigen-slice (n-500 .. n-100)
%   (mid-low variance); w_2 = w* + 100 x unit direction in the bottom-100
%   eigenvectors (direction = projection of w* onto the slice, normalized).
%   (Docstring says alpha=1000 for w_2; code uses ALPHA2=100.)
%   Accentuation: closed form I = I_0 + lam w, lam = (t - w.I_0)/||w||^2
%   (natural seed, not null seed); 13 targets linspace(mu-3 sd, mu+3 sd) of
%   f* on natural patches; 20 seeds (seed 7), 8 shown; Fig. 1B shows every
%   third target of the seed with f*(I_0) = -31.367.
%   R^2 in panels: 1 - Var(f_x - f_y)/Var(f_y) (offset-free); Fig. 1C values
%   (0.9969, 0.0481) do not match exp2_cross_eval (0.9966, -9.94): the
%   composite likely used a variant denominator -- check before final.
\todo{Write: data, constructed students, closed-form accentuation, targets,
peer-review panels.}

% ---------------------------------------------------------------------------
\subsection{Two-dimensional geometry illustration}
\label{sec:app-methods-geometry}
% Used by: Fig.~\ref{fig:geometry}, Fig.~\ref{fig:supp-ridge-path-geometry}.
% Content: Sigma = diag(1, eps=0.04); teacher; fixed design X; exact
% conditional Gaussian of beta_hat | X (mean path, 68% ellipses);
% lambda = 0.05 (kappa = 0.0518) with sigma^2 on a log grid; sigma^2 = 0.5 with
% lambda on a log grid; iso-sets of Prop.~\ref{prop:mr-iso-geometry}.
% Source: notebooks/ridge_paths_geometry_explorer.ipynb,
% tables/ridge_paths_iso_error_geometry.csv.
\todo{Write: parameters and construction of the 2D panels.}

% ---------------------------------------------------------------------------
% Pixel-ridge landscape (Fig.~\ref{fig:pixel}); subsection with its own label
% sec:app-vanhateren-landscape-methods.
\subsection{Natural-image pixel-ridge landscape experiment}
\label{sec:app-vanhateren-landscape-methods}

This subsection describes the experiment underlying Fig.~\ref{fig:pixel} and
Figs.~\ref{fig:supp-pixel-de-mc}--\ref{fig:supp-pixel-selection}.  Its
Monte Carlo component uses actual natural-image patches, rather than Gaussian
designs having the same covariance.

\paragraph{Images, teacher, and population reference.}
Images, pools, covariance and disk teacher are those of
App.~\ref{sec:app-methods-images}: a 2,000-patch training pool, a disjoint
20,000-patch population pool that defines $\Sigma$ and is used to evaluate
every fitted model, and a disk teacher with
$S=\bstar^\top\Sigma\bstar=1736.34$.

\paragraph{Noise--regularization grid and deterministic equivalents.}
We evaluated 114 response-noise levels, comprising zero and a logarithmic grid
of $\sigma^2/S\in[10^{-8},10]$, and 478 ridge penalties.  We use the
scikit-learn convention
\begin{equation}
 \widehat\beta_{\alpha,\sigma}
 =X^\top(XX^\top+\alpha I_n)^{-1}\bm y_\sigma,
 \qquad \alpha=n\lambda,
\end{equation}
with $\alpha\in[10^{-5},10^7]$ ($\lambda\in[10^{-8},10^4]$).  The effective
penalty $\kappa$ shown in Fig.~\ref{fig:pixel} is the positive solution of
$n\lambda=\kappa\{n-\mathrm{df}_1(\kappa)\}$.  At each of the 54,492 grid
points we evaluated the deterministic equivalents using the eigenvalues of
the empirical population covariance and the disk teacher's coordinates in
that eigenbasis.  The prediction-optimal and control-optimal paths minimize,
respectively, the DE prediction error and the leading-DE accentuation error
over the finite penalty grid at each noise level.  All heatmaps and slices
are direct evaluations on this grid; no smoothing or interpolation was used.

\paragraph{Natural-image Monte Carlo.}
We ran 100 paired trials.  On trial $t$, using seed $20260918+t$, we sampled
$n=1,000$ patches without replacement from the 2,000-patch training pool,
centered every pixel using that sample, and drew a centered
$\epsilon\sim\mathcal N(0,I_n)$.  For every noise level we used the response
\begin{equation}
  \bm y_\sigma=X\bstar+\sigma\epsilon .
  \label{eq:app-vh-response}
\end{equation}
Thus the design and unit-noise direction were shared across the entire noise
axis within a trial.  This pairing reduces differences caused by resampling
and makes each trial a coherent path as $\sigma$ varies.

Because $n<d$, we diagonalized the $n\times n$ kernel matrix $XX^\top$ once
per trial.  For every penalty the fitted weight is exactly affine in the
noise scale,
\begin{equation}
 \widehat\beta_{\alpha,\sigma}
 =X^\top(XX^\top+\alpha I_n)^{-1}X\bstar
 +\sigma X^\top(XX^\top+\alpha I_n)^{-1}\epsilon
 =w_{0,\alpha}+\sigma w_{1,\alpha}.
 \label{eq:app-vh-affine-noise}
\end{equation}
We therefore cached the linear or quadratic coefficients in $\sigma$ of all
required quadratic forms and reconstructed the complete noise axis exactly.
This avoids refitting 54,492 models per trial and is not an interpolation.

Let $P\in\mathbb R^{20,000\times d}$ be the centered population-patch matrix,
$N=\bstar^\top\widehat\beta$, $D=\|\widehat\beta\|^2$,
$Q=\|P\widehat\beta\|^2/20{,}000$, and
$C=(P\bstar)^\top P\widehat\beta/20{,}000$.  We report
\begin{align}
 \frac{E_{gen}}S
 &=\frac{\|P(\widehat\beta-\bstar)\|^2}{20{,}000S},
 &R^2_{gen}&=1-\frac{E_{gen}}S,
 &\operatorname{slope}_{gen}&=\frac CQ,\\
 \frac{E_{acc}}S
 &=\left(1-\frac ND\right)^2,
 &R^2_{acc}&=1-\left(\frac DN-1\right)^2,
 &\operatorname{slope}_{acc}&=\frac ND.
 \label{eq:app-vh-metrics}
\end{align}
The fitted intercept is accounted for by sample centering and by the
cross-validation leverage below, but it is omitted from these population
weight-response metrics.

\paragraph{Empirical RidgeCV.}
For each trial and noise level we selected $\alpha$ by exact leave-one-out
cross-validation over the same 478-point grid.  If $XX^\top=U\operatorname{diag}(q_j)U^\top$,
the diagonal of the fitted-intercept ridge smoother is
\begin{equation}
 h_i(\alpha)=\sum_jU_{ij}^2\frac{q_j}{q_j+\alpha}+\frac1n,
\end{equation}
and the LOOCV loss is the mean of
$[(y_i-\widehat y_i)/(1-h_i)]^2$.  As in
Eq.~\ref{eq:app-vh-affine-noise}, this loss is quadratic in $\sigma$, so its
full landscape was reconstructed from three cached coefficients.  The green
path in Fig.~\ref{fig:pixel} is the median selected $\kappa$ across trials.
Metrics labeled empirical RidgeCV were instead evaluated at each trial's own
selected penalty before aggregation; they are not evaluations at the median
penalty.  We verified the selected penalties and minimum losses against
\texttt{sklearn.linear\_model.RidgeCV} at three noise levels.

\paragraph{Aggregation and checks.}
At every grid point we saved all 100 trial values and computed their
arithmetic mean, standard error, median, and 10th and 90th percentiles.  Bands
in Fig.~\ref{fig:pixel} show the 10--90\% trial range, not a confidence
interval; markers show the arithmetic mean.  We also checked the cached
noise-polynomial evaluation against an independent direct ridge solve for
all reported metrics.  Computation used float64 linear algebra on a GPU, and
plot-ready summaries were cached separately from the rendered figure.


% ---------------------------------------------------------------------------
\subsection{Linear feature maps and prediction-null rotations}
\label{sec:app-methods-feature}
% Used by: Fig.~\ref{fig:feature}B--D, App.~\ref{sec:app-ext-gauge}.
% Theory of the construction: App.~\ref{sec:app-gauge-fragility}
% (whitened coordinates, right rotations with a fixed point,
% Cor.~\ref{cor:mr-rotation}).
% Content: base map = top-p population-PC projector (p = 500 main, 100 alt.);
% n = 1000; rowwise Givens rotations from high-variance to tail PCs, tail
% loading t = sin^2 theta; random family: 300 randomized rowwise rotations and
% pairings; noise sigma^2/S in {0.01, 0.1, 1, 10}; DE-selected alpha vs.\
% empirical CV; MC trials; computation of T_F.
% Sources: notebooks/README_top_pc_rotation.md, top_pc_rotation_experiment.ipynb,
% tables/vanhateren_top500_*, tables/vanhateren_null_*.
\todo{Write: feature maps, rotation schedule, random family, fitting and
evaluation.}

% ---------------------------------------------------------------------------
\subsection{Neuronal data}
\label{sec:app-methods-neural-data}
% Used by: Fig.~\ref{fig:phenomena} (top), Sec.~\ref{sec:pixel} (peer
% statistics), Fig.~\ref{fig:nonlinear-biological-validation},
% App.~\ref{sec:app-ext-peer}, \ref{sec:app-ext-nonlinear}.
% Sources: upstream jacob-prince/parametric-neural-control (branch jacob,
% SHA 9b6fd95; local copy $STORE_DIR/Projects/AccentuationPredRMT/
% nonlinear_control/biological_validation_v1/upstream), pnc/preproc/*.py,
% scripts/preprocessing/run_preproc.py; downstream: AccentuationPredRMT/
% explorations/nonlinear_control/validate_biology.py,
% build_biological_synopsis.py.  Data NOT from BrainScore: lab storage /
% Zenodo (record id not yet minted; CITATION.cff DOIs are placeholders).
% Subjects/areas (Monkey: recording, area, channels):
%   R red   chronic microwire array (64ch)          aIT  0,2,9,15,19
%   P paul  chronic floating microelectrode (64ch)  cIT  0,8,24,40,47
%   V venus Neuropixels (383ch)  V3/V4  9(V3),79(V4),151(V4),331(V3),355(V3)
%   L leap  Neuropixels          STS    81,282,286,306,342
%   T three0 Neuropixels         STS    56,74,120,168,204
%   Sites chosen by split-half (Spearman-Brown) reliability; R,P = exact
%   top-5; V,L,T hand-picked among reliable (rule undocumented).
%   Unit type (MUA vs sorted) not determined from code.
% Encoding (calibration) set: 969 images = 649 NSD shared1000 (COCO) +
%   259 segmented objects/animals + 61 fLoc; split 774/195
%   (encoding_stimuli_split_seed0.csv); R: 282 px @35.31 px/deg (~8 deg),
%   167 ms on / 200 ms off, (0,2) deg; ~3k-16k trials per monkey.
% Control sessions: 4-12 May 2025; 10 NSD seed images x 11 target levels per
%   site-model pair (27,720 synthesized; STS monkeys saw 43%/55%);
%   mean repeats R 3.97, P 2.29, V 4.22, L 1.22, T 1.34; calibration anchors
%   re-shown (100; 49 L; 46 T).  Target levels: linspace(q01-0.25bw,
%   q99+0.5bw, 11) of MEASURED z-scored encoding responses.
% Preprocessing (upstream): peak-window average (R,P 100-400 ms; V,L,T
%   80-250 ms); control-phase outlier rejection (|x-med|/(1.4826 MAD)>15);
%   anchorDay z-scoring (fallback to all trials if <5 anchor trials);
%   repeat averaging; Leap duplicate merge (pixel r>=0.90); firing floor clamp
%   of predictions at comparison time; NSD-style noise ceiling (undefined L,T).
% Downstream (ours, not upstream): per-site cross-session affine map fitted on
%   encoding-train anchors, frozen and shared across models; endpoints
%   control slope (OLS measured on predicted), control MSE (no refit), E_val
%   (held-out anchor MSE).  Upstream metrics: control_r, control_slope
%   (linregress), control_R2 (identity), encoding_r.
% Feature accentuation (for completeness; cite Hamblin et al. 2024, MACO):
%   Fourier parameterization of the seed (1024 px), 1/f^decay preconditioning,
%   NAdam lr 12, <=6000 steps, 8 random crops + noise per step, closest-to-
%   target stopping, (noise, decay) ladder chosen per site-model.
\todo{Write.}

% ---------------------------------------------------------------------------
\subsection{Encoding models}
\label{sec:app-methods-models}
% Content: Table~\ref{tab:app-models}.  Loaders:
% Closed-loop-visual-insilico/core/model_load_utils.py:load_model_transform.
% Preprocessing: resize 224 + ImageNet norm (RN50, robust RN50, DINO RN50,
% DINOv2, AlexNet); RADIO resize only; CLIP-RN50, CLIPAG, SigLIP2, RegNetY:
% library eval transforms (NB: resnet50_clip used ImageNet norm at synthesis
% but CLIP transform at fitting -- upstream inconsistency, footnote?).
% Fitting: flattened layer activations -> PCA 750 (fit on 774 train images,
% centered, not whitened) -> sklearn RidgeCV(alphas = 14 log-spaced in
% [1e-4, 1e9], efficient LOO, alpha per channel); layer chosen per channel by
% held-out (195-image) R^2 -- the same split later reports encoding r.
% Candidate layers: Bottleneck blocks (ResNets, RegNet), ReLU/Conv/MaxPool
% (AlexNet), transformer blocks (ViTs; + attention pool for SigLIP2).
\todo{Write; fill Table~\ref{tab:app-models}.}

\begin{table}[ht]
\centering\small
\begin{tabular}{@{}p{0.2\linewidth}p{0.22\linewidth}p{0.36\linewidth}c@{}}
\toprule
Model & Architecture & Training & \# geometries \\
\midrule
AlexNet (untrained) & AlexNet (TF-slim variant) & random initialization$^\dagger$ & 11\\
ResNet-50 & ResNet-50 & ImageNet supervised & 9\\
Robust ResNet-50 & ResNet-50 & ImageNet adversarial, $\ell_\infty$, $\epsilon=8/255$ & 6\\
CLIP RN50 & ResNet-50 & CLIP contrastive & 9\\
DINO RN50 & ResNet-50 & DINO self-supervised & 6\\
RegNetY-640 & RegNetY-640 & SEER self-supervised & 9\\
CLIPAG ViT-B/32 & ViT-B/32 & adversarially fine-tuned CLIP & 7\\
DINOv2 ViT-B/14-reg & ViT-B/14 + registers & DINOv2 self-supervised & 9\\
SigLIP2 ViT-B/16 & ViT-B/16 & SigLIP2 & 10\\
RADIO v2.5-B & ViT-B & RADIO distillation & 8\\
\bottomrule
\end{tabular}
\caption{\textbf{Encoding models.} \# geometries: distinct (layer, PCA)
feature maps across the 25 sites (84 in total). $^\dagger$Checkpoint
\texttt{AlexNet\_training\_seed\_01} (Brain-Score model-variability weights)
is at its random initialization. \todo{citations; checkpoint sources of the
robust ResNet-50 and CLIPAG}}
\label{tab:app-models}
\end{table}

% ---------------------------------------------------------------------------
\subsection{Computing the gradient-geometry diagnostic}
\label{sec:app-methods-nonlinear}
% Used by: Fig.~\ref{fig:nonlinear-biological-validation},
% App.~\ref{sec:app-ext-nonlinear}.  Formulas and estimator definitions:
% App.~\ref{sec:app-nonlinear-diagnostic} (Eqs.~\ref{eq:nonlinear-local-fragility},
% \ref{eq:nonlinear-neighborhood-estimators}, \ref{eq:nonlinear-stein-estimator}).
% Content: kappa from the CV alpha and the training PCA spectrum
% (kappa(1 - df_1/n) = alpha/n, n = 774); one VJP per PC (750, batched);
% 224x224 RGB grid in [0,1] units (d = 150,528); average over the 10 seeds;
% neighborhood scales tau in {0.5, 2, 8, 16}/255, R = 64 directions,
% L = 8 centers; debiased U-statistics; finite-difference check; antithetic
% Stein with R = 512; E_val; V_phi = E_val T / n.  Checks: 84 unique
% model--layer geometries, residual <= 6e-8; compute ~1.4 H100-hours.
\todo{Write.}

% ---------------------------------------------------------------------------
\subsection{Neuronal peer-review analysis}
\label{sec:app-methods-peer}
% Used by: Sec.~\ref{sec:phenomena}, peer paragraph of Sec.~\ref{sec:pixel},
% Fig.~\ref{fig:supp-neural-peer-matrix}, Fig.~\ref{fig:supp-neural-self-vs-peer}.
% Content: 25 sites x 10 generators x 10 reviewers; reviewer prediction on
% the generator's stimuli through the same frozen affine map; MSE, r,
% identity R^2, slope; equal-weight site means; self (diagonal) vs.\ peer
% tests: reviewer-label permutation, Wilcoxon over generators, subject
% sign-flip, subject bootstrap.
\todo{Write.}

% ---------------------------------------------------------------------------
\subsection{Statistical analysis}
\label{sec:app-methods-stats}
% Content: log predictors; within-site centering of predictor and outcome;
% Pearson / Spearman / site-clustered OLS; model subsets 10/9/8/7;
% Benjamini--Hochberg within endpoint x subset; image bootstrap;
% leave-one-monkey-out / leave-one-model-out; repeated-measures caveat.
\todo{Write.}

% ---------------------------------------------------------------------------
\subsection{Software and compute}
\label{sec:app-methods-compute}
\todo{Write (short).}
